\section{Ablation of the learned residue-view count}
\label{app:v4-k-ablation}

\paragraph{Motivation and capacity.}
The hyperparameter $K$ determines how many learned residue summaries are
available before view fusion. The global mean preserves broad sequence
context, and SLEEC provides a fixed functional-residue prior; the learned
queries can adapt their pooling to enzyme--reaction associations. Multiple
queries can preserve complementary patterns that a single weighted mean
would merge. This is a capacity motivation: neither disjoint residue sets
nor named biochemical roles are imposed, and additional views may be redundant.

Each extra view adds a 256-dimensional query, an independent projection,
and wider gate input and output layers. At the stated hidden and output
widths, the trainable multiview enzyme module has
$P(K)=3{,}159{,}298+395{,}008K$ parameters: 3.554, 3.949, 4.739, and
6.319 million for $K=1,2,4,8$, respectively. These counts exclude frozen
feature resources and SLEEC, the reaction encoder, and the phase-2 heads.
Thus $K=4$ permits several adaptive summaries at a lower parameter cost than
$K=8$, while keeping the final neural embedding dimension fixed at 512.

\paragraph{Matched protocol.}
We train $K\in\{1,2,4,8\}$ separately on EnzymeMap and the three ReactZyme
splits, giving 16 fresh fits with seed 42. The $K=4$ reference is retrained
within this study. All arms use the proposed association-only objective.
Within each target, only $K$ and its induced
projection and gate dimensions vary. Frozen inputs, association graph,
1,022-residue input limit, SLEEC/global views, losses, and training and
inference settings are held fixed: phase 1 uses 10 epochs and batch 512;
phase 2 uses 100 updates; inference uses $(\nu,c,\alpha)=(3,0.5,0.4)$
(Appendix~\ref{app:v4-recipe}). We report every arm at those fixed checkpoints
on the same complete benchmark candidate pools. These fresh-run ablation
values are distinct from the historical checkpoints in the primary
benchmark comparison.

\begin{table}[htbp]
\centering
\caption{Learned residue-view count on ReactZyme. Entries are query-averaged
all-positive test MRR; $R\!\to\!E$ and $E\!\to\!R$ denote the retrieval
directions. All rows use seed 42 and fresh target-specific training.
Bold marks the largest value in each column, without implying significance.}
\label{tab:v4-k-reactzyme}
\begin{tabular}{crrrrrr}
\toprule
 & \multicolumn{2}{c}{Reaction-Sim} & \multicolumn{2}{c}{Enzyme-Sim} & \multicolumn{2}{c}{Time} \\
\cmidrule(lr){2-3}\cmidrule(lr){4-5}\cmidrule(lr){6-7}
$K$ & $R\!\to\!E$ & $E\!\to\!R$ & $R\!\to\!E$ & $E\!\to\!R$ & $R\!\to\!E$ & $E\!\to\!R$ \\
\midrule
1 & 0.4000 & \textbf{0.5419} & 0.6632 & 0.9652 & 0.5374 & \textbf{0.8042} \\
2 & \textbf{0.4045} & 0.5158 & 0.6672 & 0.9626 & 0.5381 & 0.7975 \\
4 & 0.3947 & 0.5410 & 0.6662 & \textbf{0.9704} & 0.5383 & 0.7897 \\
8 & 0.3926 & 0.5010 & \textbf{0.6674} & 0.9613 & \textbf{0.5398} & 0.7911 \\
\bottomrule
\end{tabular}
\end{table}

\begin{table}[htbp]
\centering
\caption{Learned residue-view count on EnzymeMap. BEDROC values are
percentages; EF5 and EF10 are fold enrichments. The full library contains
1,521 queries and 261,907 candidates; exclusion of training identifiers
leaves 1,337 queries and 252,113 candidates. Bold marks the largest value
within each setting and metric; all rows use seed 42.}
\label{tab:v4-k-enzymemap}
\begin{tabular}{crrrr}
\toprule
$K$ & BEDROC$_{85}$ & BEDROC$_{20}$ & EF5 & EF10 \\
\midrule
\multicolumn{5}{l}{\emph{Full library}} \\
1 & 52.22 & 70.78 & 15.82 & 8.67 \\
2 & \textbf{57.97} & \textbf{75.61} & 16.73 & 8.90 \\
4 & 57.23 & 75.54 & \textbf{16.86} & \textbf{8.96} \\
8 & 50.40 & 69.05 & 15.59 & 8.49 \\
\midrule
\multicolumn{5}{l}{\emph{Training-ID-excluded library}} \\
1 & 47.03 & 67.09 & 15.22 & 8.45 \\
2 & \textbf{53.64} & 72.61 & 16.23 & 8.72 \\
4 & 52.87 & \textbf{72.63} & \textbf{16.41} & \textbf{8.79} \\
8 & 44.97 & 65.15 & 14.90 & 8.23 \\
\bottomrule
\end{tabular}
\end{table}

\paragraph{Results and retained default.}
Tables~\ref{tab:v4-k-reactzyme} and~\ref{tab:v4-k-enzymemap} show no uniformly
best count. On EnzymeMap, $K=2$ gives the highest BEDROC$_{85}$ in both
settings, whereas $K=4$ gives the highest EF5 and EF10. Relative to $K=2$,
$K=4$ loses 0.74 and 0.78 BEDROC$_{85}$ percentage points in the full and
excluded libraries, respectively. On ReactZyme, $K=4$ improves
$E\!\to\!R$ MRR over $K=2$ on Reaction-Sim and Enzyme-Sim, but reduces
Reaction-Sim $R\!\to\!E$ and Time $E\!\to\!R$ MRR. Increasing to $K=8$
yields small gains in Enzyme-Sim and Time $R\!\to\!E$ MRR, but lower
EnzymeMap scores and lower Reaction-Sim MRR in both directions than $K=4$.
We therefore retain the original $K=4$ default as a practical compromise
across screening and bidirectional retrieval, rather than claim that it
maximizes every metric.

\paragraph{Scope of the evidence.}
This is a single-seed sensitivity study, not an estimate of robustness or
statistical significance. Increasing $K$ also increases parameter count;
the comparison is not parameter-matched, and a common seed does not ensure
identical initialization across different shapes. Since every arm retains
at least one learned view, this experiment does not isolate the benefit of
adding learned views relative to a $K=0$ model. The interpretation applies
to the proposed association-only training setting; it does not establish
an optimal count across random seeds or unrelated downstream tasks.
